\subsection{Polar and Singular Value Decompositions}

Throughout this subsection, $\mathbb{F}\in\{\mathbb{R},\mathbb{C}\}$, and all matrix dimensions are positive integers.\\
Every inner-product space is finite-dimensional over $\mathbb{F}$.

\begin{theorem} \label{thm:polar-decomposition} \hlt{Polar Decomposition}\\
Let $T\in\LL(\mathcal{V})$. Then $(\exists S\in\LL(\mathcal{V}))\,[S\text{ is an isometry}\land T=S\sqrt{T^*T}]$.
\begin{proof}
\begin{enumerate}[label=(\roman*)]
\setlength{\itemsep}{0pt}
\item \hlt{Positive Factor}: Set $P\defeq\sqrt{T^*T}$ by Theorem \ref{thm:positive-operator-characterization}(v)$\Rightarrow$(i) and Proposition \ref{prop:positive-square-root-unique}.\\
$(\forall\mathbf{v}\in\mathcal{V})\,\norm{T\mathbf{v}}^2=\innerprod{P^2\mathbf{v},\mathbf{v}}=\norm{P\mathbf{v}}^2$ by $P=P^*$. $\ker(P)=\ker(T)$ by this norm identity.
\item \hlt{Range Isometry}: The rule $S_1(P\mathbf{v})\defeq T\mathbf{v}$ defines a linear bijection $\operatorname{range}(P)\to\operatorname{range}(T)$ by (i).\\
$S_1$ is an isometry by $\norm{S_1(P\mathbf{v})}=\norm{P\mathbf{v}}$ from (i).
\item \hlt{Complement Extension}: $\dim\operatorname{range}(P)^\perp=\dim\operatorname{range}(T)^\perp$ by (i), Theorem \ref{thm:rank-nullity}, and Corollary \ref{cor:dimension-of-orthogonal-complement}.\\
Set $d\defeq\dim\operatorname{range}(P)^\perp$. Choose orthonormal bases $(\mathbf{a}_j)_{j=1}^{d}$ and $(\mathbf{b}_j)_{j=1}^{d}$ of the two complements by Corollary \ref{cor:finite-dimensional-inner-product-space-has-orthonormal-basis}.\\
Define $S_0\mathbf{a}_j\defeq\mathbf{b}_j$. Then $(\forall i,j\in J_d)\,\innerprod{S_0\mathbf{a}_i,S_0\mathbf{a}_j}=\innerprod{\mathbf{a}_i,\mathbf{a}_j}=\delta_{ij}$. Thus $S_0$ is an isometry by sesquilinearity and Theorem \ref{thm:isometry-iff-inner-product-preserving}.
\item \hlt{Global Isometry}: $\mathcal{V}=\operatorname{range}(P)\odot\operatorname{range}(P)^\perp$ by Theorem \ref{thm:orthogonal-direct-sum-decomposition}. Define $S(\mathbf{x}+\mathbf{y})\defeq S_1\mathbf{x}+S_0\mathbf{y}$.\\
$S$ is an isometry by the orthogonal source and target decompositions. Finally, $(\forall\mathbf{v}\in\mathcal{V})\,SP\mathbf{v}=S_1(P\mathbf{v})=T\mathbf{v}$.
\end{enumerate}
\hlt{Structural Commentary}: The Gram norm identity builds a range isometry, extended across equidimensional complements.\\
\hlt{Key Insight}: $\ker\sqrt{T^*T}=\ker T$ synchronizes the two range dimensions.
\end{proof}
\end{theorem}

\begin{corollary} \label{cor:operator-polar-decomposition-uniqueness} \hlt{Uniqueness in Operator Polar Decomposition}\\
Let $T\in\LL(\mathcal{V})$, and set $\abs{T}\defeq\sqrt{T^*T}$ and $\abs{T^*}\defeq\sqrt{TT^*}$.
\begin{enumerate}[label=(\roman*)]
\setlength{\itemsep}{0pt}
\item If $T=SP$, $S^*S=I_{\mathcal{V}}$, and $P\geq0$, then $P=\abs{T}$; hence the positive right polar factor is unique.
\item $(\exists R\in\LL(\mathcal{V}))\,[R\text{ is an isometry}\land T=\abs{T^*}R]$, and $\abs{T^*}$ is the unique positive left polar factor.
\item If $T$ is invertible, then both isometric factors are unique: $S=T\abs{T}^{-1}$ and $R=\abs{T^*}^{-1}T$.
\end{enumerate}
\begin{proof}
\begin{enumerate}[label=(\roman*)]
\setlength{\itemsep}{0pt}
\item If $T=SP$, then $T^*T=PS^*SP=P^2$, so $P=\abs{T}$ by Proposition \ref{prop:positive-square-root-unique}.
\item Fix $T^*=S_0\abs{T^*}$ by Theorem \ref{thm:polar-decomposition}.\\
Then $R\defeq S_0^*$ is an isometry and $T=\abs{T^*}R$ by Theorem \ref{thm:finite-dimensional-isometry-iff-adjoint-product-identity} and Proposition \ref{prop:basic-isometry-properties}(iii)--(iv).\\
If $T=Q\widetilde R$ with $Q\geq0$ and $\widetilde R$ isometric, then $TT^*=Q^2$ by $\widetilde R\widetilde R^*=I_{\mathcal{V}}$; hence $Q=\abs{T^*}$ by Proposition \ref{prop:positive-square-root-unique}.
\item If $T$ is invertible, then $\ker\abs{T}=\ker\abs{T^*}=\{\mathbf{0}\}$ by invertibility of their squares.\\
Thus both moduli are invertible by Theorem \ref{thm:finite-dimensional-operator-invertible-equivalences}, and $S=T\abs{T}^{-1}$, $R=\abs{T^*}^{-1}T$.
\end{enumerate}
\hlt{Structural Commentary}: Gram squaring fixes each modulus; invertibility then isolates its isometry.\\
\hlt{Key Insight}: Positivity selects one square root of each Gram operator.
\end{proof}
\end{corollary}

Figure \ref{fig:polar-decomposition-operator} summarizes the range isometry used to separate the positive stretch from the isometric factor.

\begin{illustrationfigure}[H]
  \centering
  \begin{tikzpicture}[
  >=Stealth,
  x=0.86cm,
  every node/.style={font=\small},
  space/.style={draw=gray!55, rounded corners=3pt, minimum width=2.2cm, minimum height=0.9cm, align=center},
  map/.style={->, line width=0.9pt, blue!70!black},
  iso/.style={->, line width=0.9pt, teal!70!black}
]
  \node[space] (V) at (0,0) {$\mathcal{V}$};
  \node[space] (P) at (3.3,0) {$\operatorname{range}\sqrt{T^*T}$};
  \node[space] (T) at (6.9,0) {$\operatorname{range}T$};

  \draw[map] (V) -- (P) node[midway, above=11pt, fill=white, inner sep=1pt] {$\sqrt{T^*T}$};
  \draw[iso] (P) -- (T) node[midway, above] {$S_1$};
  \draw[map, out=-45, in=-135, looseness=0.85] (V.south east)
    to node[pos=0.5, above, fill=white, inner sep=1pt] {$T$} (T.south west);
\end{tikzpicture}

  \caption{Polar decomposition builds $S_1$ using $\norm{T\mathbf{v}}=\norm{\sqrt{T^*T}\mathbf{v}}$.}
  \label{fig:polar-decomposition-operator}
\end{illustrationfigure}

\begin{definition} \label{def:singular-values} \hlt{Singular Values}\\
For $T\in\LL(\mathcal{U},\mathcal{V})$, the \hlt{singular values} of $T$ are the eigenvalues of $\sqrt{T^*T}\in\LL(\mathcal{U})$, listed with their eigenspace multiplicities.
\end{definition}

\begin{theorem} \label{thm:singular-value-decomposition} \hlt{Singular Value Decomposition}\\
Let $T\in\LL(\mathcal{V})$, set $n\defeq\dim\mathcal{V}$, and list its singular values as $s_1,\ldots,s_n$.\\
Then $\exists$ orthonormal bases $(\mathbf{e}_1,\ldots,\mathbf{e}_n)$ and $(\mathbf{f}_1,\ldots,\mathbf{f}_n)$ of $\mathcal{V}$ such that $(\forall\mathbf{v}\in\mathcal{V})\,T\mathbf{v}=\sum\limits_{j=1}^{n}s_j\innerprod{\mathbf{v},\mathbf{e}_j}\mathbf{f}_j$.
\begin{proof}
Set $\abs{T}\defeq\sqrt{T^*T}$. Fix an orthonormal eigenbasis $(\mathbf{e}_1,\ldots,\mathbf{e}_n)$ of $\abs{T}$ with $(\forall j\in J_n)\,\abs{T}\mathbf{e}_j=s_j\mathbf{e}_j$ by Theorem \ref{thm:real-spectral-theorem} when $\mathbb{F}=\mathbb{R}$ and Theorem \ref{thm:complex-spectral-theorem} when $\mathbb{F}=\mathbb{C}$. Fix $T=S\abs{T}$ by Theorem \ref{thm:polar-decomposition}, and set $(\forall j\in J_n)\,\mathbf{f}_j\defeq S\mathbf{e}_j$.\\
The list $(\mathbf{f}_1,\ldots,\mathbf{f}_n)$ is an orthonormal basis by Theorem \ref{thm:isometry-iff-inner-product-preserving} and Proposition \ref{prop:basic-isometry-properties}(iii). $(\forall\mathbf{v}\in\mathcal{V})\,T\mathbf{v}=\sum\limits_{j=1}^{n}s_j\innerprod{\mathbf{v},\mathbf{e}_j}\mathbf{f}_j$ by Theorem \ref{thm:orthonormal-basis-coordinates}.\\
\hlt{Structural Commentary}: Spectral coordinates diagonalise $\abs{T}$; the polar isometry transports its eigenbasis.\\
\hlt{Key Insight}: $T\mathbf{e}_j=s_j\mathbf{f}_j$ pairs stretch with direction.
\end{proof}
\end{theorem}

Figure \ref{fig:singular-value-decomposition} shows each right singular direction stretched and carried isometrically to its left singular direction.

\begin{illustrationfigure}[H]
  \centering
  \begin{tikzpicture}[
  >=Stealth,
  x=0.86cm,
  every node/.style={font=\small},
  vec/.style={draw=gray!55, rounded corners=3pt, inner sep=4pt, align=center},
  stretch/.style={->, line width=0.9pt, blue!70!black},
  isoarrow/.style={->, line width=0.9pt, teal!70!black}
]
  \node[vec] (e) at (0,0) {right singular\\direction $\mathbf{e}_j$};
  \node[vec] (se) at (3.6,0) {stretch\\$s_j\mathbf{e}_j$};
  \node[vec] (f) at (7.2,0) {left singular\\direction $s_j\mathbf{f}_j$};

  \draw[stretch] (e) -- (se) node[midway, above] {$\sqrt{T^*T}$};
  \draw[isoarrow] (se) -- (f) node[midway, above] {$S$};
\end{tikzpicture}

  \caption{Since $T=S\sqrt{T^*T}$, each $\mathbf{e}_j$ maps to $s_j\mathbf{f}_j$.}
  \label{fig:singular-value-decomposition}
\end{illustrationfigure}

\begin{proposition} \label{prop:singular-values-from-t-star-t} \hlt{Singular Values from $T^*T$}\\
The singular values of $T\in\LL(\mathcal{U},\mathcal{V})$ are the nonnegative square roots of the eigenvalues of $T^*T$, with the same eigenspace multiplicities.
\begin{proof}
Set $n\defeq\dim\mathcal{U}$ and $G\defeq T^*T$. Then $G\geq0$ by Theorem \ref{thm:positive-operator-characterization}(v)$\Rightarrow$(i).\\
Fix an orthonormal eigenbasis $(\mathbf{e}_j)_{j=1}^{n}$ with $G\mathbf{e}_j=\lambda_j\mathbf{e}_j$ by Theorem \ref{thm:real-spectral-theorem} when $\mathbb{F}=\mathbb{R}$ and Theorem \ref{thm:complex-spectral-theorem} when $\mathbb{F}=\mathbb{C}$.\\
$(\forall j\in J_n)\,\lambda_j\geq0$ by Theorem \ref{thm:positive-operator-characterization}(i)$\Rightarrow$(ii). Define $(\forall j\in J_n)\,R\mathbf{e}_j\defeq\sqrt{\lambda_j}\mathbf{e}_j$.\\
Then $R\geq0$ by Theorem \ref{thm:positive-operator-characterization}(ii)$\Rightarrow$(i), and $R^2=G$; hence $R=\sqrt{G}$ by Proposition \ref{prop:positive-square-root-unique}.\\
The eigenspace multiplicities are preserved by $(\forall\lambda\geq0)\,E_{\lambda}(G)=E_{\sqrt{\lambda}}(R)$.\\
\hlt{Structural Commentary}: One Gram eigenbasis constructs the root and identifies every eigenspace.\\
\hlt{Key Insight}: $\lambda\mapsto\sqrt{\lambda}$ is injective on the nonnegative spectrum.
\end{proof}
\end{proposition}

\begin{theorem} \label{thm:rectangular-operator-singular-vector-decomposition} \hlt{Singular Vectors for a Linear Map}\\
Let $T\in\LL(\mathcal{U},\mathcal{V})$ have rank $r$. Set $n\defeq\dim\mathcal{U}$, $m\defeq\dim\mathcal{V}$, $\mathcal{U}_+\defeq\operatorname{range}(T^*)$, $\mathcal{U}_0\defeq\ker(T)$, $\mathcal{V}_+\defeq\operatorname{range}(T)$, and $\mathcal{V}_0\defeq\ker(T^*)$.\\
Then $\exists$ ordered orthonormal bases $\beta=(\mathbf{u}_1,\ldots,\mathbf{u}_n)$ and $\gamma=(\mathbf{v}_1,\ldots,\mathbf{v}_m)$\\
with $\operatorname{span}\{\mathbf{u}_j:j\in J_r\}=\mathcal{U}_+$, $\operatorname{span}\{\mathbf{u}_j:j\in J_n\setminus J_r\}=\mathcal{U}_0$, $\operatorname{span}\{\mathbf{v}_j:j\in J_r\}=\mathcal{V}_+$, and $\operatorname{span}\{\mathbf{v}_j:j\in J_m\setminus J_r\}=\mathcal{V}_0$,\\
and positive scalars $(s_k)_{k\in J_r}$ such that $(\forall k\in J_r)\,T\mathbf{u}_k=s_k\mathbf{v}_k$ and $T^*\mathbf{v}_k=s_k\mathbf{u}_k$.\\
The $s_k$ are the positive singular values; $\mathbf{u}_k$ and $\mathbf{v}_k$ are right and left singular vectors.
\begin{proof}
\begin{enumerate}[label=(\roman*)]
\setlength{\itemsep}{0pt}
\item \hlt{Gram Reduction}: For $G\defeq T^*T$, $G\geq0$, $\ker(G)=\ker(T)$ by Theorem \ref{thm:positive-operator-characterization}, Proposition \ref{prop:adjoint-product-kernels-ranges-and-projections}; $\operatorname{rank}(G)=r$ by Theorem \ref{thm:rank-nullity}.
\item \hlt{Right Singular Basis}: Choose orthonormal eigenbasis $(\mathbf{u}_1,\ldots,\mathbf{u}_n)$ of $G$ by Theorem \ref{thm:real-spectral-theorem} if $\mathbb{F}=\mathbb{R}$, Theorem \ref{thm:complex-spectral-theorem} if $\mathbb{F}=\mathbb{C}$.\\
Order it so that $G\mathbf{u}_k=s_k^2\mathbf{u}_k$ for $k\in J_r$ and $G\mathbf{u}_k=\mathbf{0}$ for $k\in J_n\setminus J_r$.\\
These groups span $\mathcal{U}_+$ and $\mathcal{U}_0$ by Proposition \ref{prop:adjoint-kernel-range-orthogonality}(iii), and the $s_k$ are the positive singular values by Proposition \ref{prop:singular-values-from-t-star-t}.
\item \hlt{Left Singular Vectors}: Set $(\forall k\in J_r)\,\mathbf{v}_k\defeq s_k^{-1}T\mathbf{u}_k$. Then $(\forall i,j\in J_r)\,\innerprod{\mathbf{v}_i,\mathbf{v}_j}=\frac{\innerprod{\mathbf{u}_i,G\mathbf{u}_j}}{s_is_j}=\delta_{ij}$.\\
$(\forall k\in J_r)\,T\mathbf{u}_k=s_k\mathbf{v}_k$ and $T^*\mathbf{v}_k=s_k\mathbf{u}_k$ by the definitions of $\mathbf{v}_k$ and $G$.
\item \hlt{Codomain Completion}: The identity $\mathcal{V}_+=T(\mathcal{U}_+)=\operatorname{span}\{\mathbf{v}_k:k\in J_r\}$ holds by $\mathcal{U}=\mathcal{U}_+\odot\mathcal{U}_0$.\\
Complete $(\mathbf{v}_1,\ldots,\mathbf{v}_r)$ orthonormally by Corollary \ref{cor:orthonormal-list-extends-to-basis}. The remaining vectors lie in $\mathcal{V}_+^\perp=\mathcal{V}_0$ by Proposition \ref{prop:adjoint-kernel-range-orthogonality}(i).
\end{enumerate}
\hlt{Structural Commentary}: The Gram spectrum splits the domain; normalized images build the active codomain basis.\\
\hlt{Key Insight}: $T^*T\mathbf{u}_k=s_k^2\mathbf{u}_k$ forces both singular-vector relations.
\end{proof}
\end{theorem}

Figure \ref{fig:singular-subspaces-adjoint-symmetry} summarizes the paired active subspaces and kernels for $T$ and $T^*$.

\begin{illustrationfigure}[H]
  \centering
  % Source note: redrawn from the singular-subspace schematic in Roman, Advanced Linear Algebra, Chapter 17.
\begin{tikzpicture}[
  >=Stealth,
  x=0.88cm,
  y=0.90cm,
  every node/.style={font=\small},
  space/.style={draw=gray!65, fill=gray!7, rounded corners=3pt, minimum width=2.22cm, minimum height=0.68cm, align=center},
  active/.style={draw=blue!60!black, fill=blue!6},
  kernel/.style={draw=gray!65, fill=gray!10},
  map/.style={->, line width=0.95pt, blue!70!black},
  back/.style={->, line width=0.95pt, teal!70!black},
  zero/.style={->, line width=0.9pt, gray!70}
]
  \node[font=\bfseries] at (0,1.55) {$\mathcal{U}$};
  \node[font=\bfseries] at (6.2,1.55) {$\mathcal{V}$};

  \node[space, active] (uactive) at (0,0.72)
    {$\operatorname{range}(T^*)$\\$\mathbf{u}_1,\ldots,\mathbf{u}_r$};
  \node[space, kernel] (uker) at (0,-0.72)
    {$\ker(T)$\\$\mathbf{u}_{r+1},\ldots,\mathbf{u}_n$};

  \node[space, active] (vactive) at (6.2,0.72)
    {$\operatorname{range}(T)$\\$\mathbf{v}_1,\ldots,\mathbf{v}_r$};
  \node[space, kernel] (vker) at (6.2,-0.72)
    {$\ker(T^*)$\\$\mathbf{v}_{r+1},\ldots,\mathbf{v}_m$};

  \draw[map] ([yshift=0.13cm]uactive.east) -- node[above] {$T\mathbf{u}_k=s_k\mathbf{v}_k$} ([yshift=0.13cm]vactive.west);
  \draw[back] ([yshift=-0.13cm]vactive.west) -- node[below] {$T^*\mathbf{v}_k=s_k\mathbf{u}_k$} ([yshift=-0.13cm]uactive.east);

  \draw[zero] (uker.east) -- node[above] {$T=0$} (vker.west);
  \draw[zero] ([yshift=-0.22cm]vker.west) -- node[below] {$T^*=0$} ([yshift=-0.22cm]uker.east);
\end{tikzpicture}

  \caption{The singular-vector bases split the domain and codomain into active singular subspaces and kernels.}
  \label{fig:singular-subspaces-adjoint-symmetry}
\end{illustrationfigure}

\begin{definition} \label{def:matrix-modulus} \hlt{Matrix Modulus}\\
Let $A\in\operatorname{Mat}(m,n,\mathbb{F})$. The matrices $A^*A$ and $AA^*$ are positive semidefinite by Proposition \ref{prop:rectangular-gram-products-are-psd}(i)--(ii).\\
Their unique positive semidefinite square roots exist by Corollary \ref{cor:positive-semidefinite-square-root-formula}(i).\\
The \hlt{moduli} of $A$ and $A^*$ are $\abs{A}\defeq(A^*A)^{1/2}\in\operatorname{Mat}(n,n,\mathbb{F})$ and $\abs{A^*}\defeq(AA^*)^{1/2}\in\operatorname{Mat}(m,m,\mathbb{F})$.
\end{definition}

\begin{definition} \label{def:matrix-singular-values} \hlt{Matrix Singular Values}\\
Let $A\in\operatorname{Mat}(m,n,\mathbb{F})$, $q\defeq\min(m,n)$, and $r\defeq\operatorname{rank}(A)$.\\
The matrix $\abs{A}$ has exactly $r$ positive eigenvalues by Theorem \ref{thm:gram-factor-kernel-rank-column-data}(iii) and Corollary \ref{cor:positive-semidefinite-square-root-formula}(ii).\\
These eigenvalues, counted with multiplicity, are the \hlt{positive singular values} of $A$.\\
For $r\geq1$, list them as $\sigma_1\geq\cdots\geq\sigma_r>0$; for $r=0$, this list is empty.\\
The full singular-value list has length $q$; if $r<q$, set $\sigma_{r+1}=\cdots=\sigma_q=0$.
\end{definition}

\begin{theorem} \label{thm:matrix-singular-value-decomposition} \hlt{Matrix Singular Value Decomposition}\\
Let $A\in\operatorname{Mat}(m,n,\mathbb{F})$ have singular values $\sigma_1\geq\cdots\geq\sigma_q\geq0$, where $q\defeq\min(m,n)$. Then $\exists$ unitary matrices $V\in\operatorname{Mat}(m,m,\mathbb{F})$ and $W\in\operatorname{Mat}(n,n,\mathbb{F})$ such that $A=V\Sigma W^*$, where the rectangular diagonal matrix $\Sigma\in\operatorname{Mat}(m,n,\mathbb{R})$ satisfies $(\forall j\in J_q)\,\Sigma_{jj}=\sigma_j$.
\begin{proof}
\begin{enumerate}[label=(\roman*)]
\setlength{\itemsep}{0pt}
\item \hlt{Gram Spectrum}: Set $r\defeq\operatorname{rank}(A)$. If $r=0$, take $(V,W,\Sigma)=(I_m,I_n,0)$; assume $r\geq1$ below.\\
Choose a unitary eigenvector matrix $W=[\mathbf{w}_1\ \cdots\ \mathbf{w}_n]$ for $A^*A$ by Theorem \ref{thm:real-spectral-theorem} when $\mathbb{F}=\mathbb{R}$ and Theorem \ref{thm:complex-spectral-theorem} when $\mathbb{F}=\mathbb{C}$.\\
Then $W^*A^*AW=\operatorname{diag}(\sigma_1^2,\ldots,\sigma_r^2,0,\ldots,0)$ by Corollary \ref{cor:positive-semidefinite-square-root-formula}(ii).
\item \hlt{Left Singular Vectors}: Set $(\forall j\in J_r)\,\mathbf{v}_j\defeq\sigma_j^{-1}A\mathbf{w}_j$. Then $(\forall i,j\in J_r)\,\mathbf{v}_i^*\mathbf{v}_j=\frac{\mathbf{w}_i^*A^*A\mathbf{w}_j}{\sigma_i\sigma_j}=\delta_{ij}$.
\item \hlt{Unitary Completion}: Complete $(\mathbf{v}_1,\ldots,\mathbf{v}_r)$ to an orthonormal basis by Corollary \ref{cor:orthonormal-list-extends-to-basis}.\\
The matrix $V$ with these columns is unitary by Theorem \ref{thm:unitary-matrix-equivalent-conditions}(ii)$\Rightarrow$(i).
\item \hlt{Factorisation}: For $j\in J_n\setminus J_r$, $\norm{A\mathbf{w}_j}_2^2=\mathbf{w}_j^*A^*A\mathbf{w}_j=0$. $AW=V\Sigma$ by (ii) and the preceding identity, so $A=V\Sigma W^*$.
\end{enumerate}
\hlt{Structural Commentary}: The Gram eigenbasis fixes $W,\Sigma$; normalized images construct $V$.\\
\hlt{Key Insight}: $\mathbf{v}_j=\sigma_j^{-1}A\mathbf{w}_j$ orthonormalizes the nonzero images.
\end{proof}
\end{theorem}

Figure \ref{fig:matrix-svd-unit-ball-to-ellipse} traces the unit ball through the three SVD factors.

\begin{illustrationfigure}[H]
  \centering
  \begin{tikzpicture}[
  >=Stealth,
  x=0.90cm,
  y=0.90cm,
  every node/.style={font=\small},
  map/.style={->, line width=0.9pt, blue!70!black},
  axis/.style={->, gray!70},
  vec/.style={->, line width=1.0pt},
  outline/.style={draw=gray!75, line width=0.8pt}
]
  \begin{scope}[shift={(0,0)}]
    \draw[axis] (-1.5,0) -- (1.65,0);
    \draw[axis] (0,-1.35) -- (0,1.5);
    \draw[outline] (0,0) circle (1);
    \draw[vec, teal!80!black] (0,0) -- (0.95,0) node[below right] {$\mathbf{w}_1$};
    \draw[vec, teal!80!black] (0,0) -- (0,0.95) node[above left] {$\mathbf{w}_2$};
    \node at (0,-1.75) {unit ball};
  \end{scope}

  \begin{scope}[shift={(4.1,0)}]
    \draw[axis] (-1.5,0) -- (1.65,0);
    \draw[axis] (0,-1.35) -- (0,1.5);
    \draw[outline] (0,0) ellipse (1.35 and 0.55);
    \draw[vec, red!75!black] (0,0) -- (1.35,0) node[below right] {$\sigma_1\mathbf{e}_1$};
    \draw[vec, red!75!black] (0,0) -- (0,0.55) node[above left] {$\sigma_2\mathbf{e}_2$};
    \node at (0,-1.75) {diagonal stretch};
  \end{scope}

  \begin{scope}[shift={(8.4,0)}, rotate=22]
    \draw[axis] (-1.65,0) -- (1.8,0);
    \draw[axis] (0,-1.35) -- (0,1.5);
    \draw[outline] (0,0) ellipse (1.35 and 0.55);
    \draw[vec, purple!75!black] (0,0) -- (1.35,0) node[below right] {$\sigma_1\mathbf{v}_1$};
    \draw[vec, purple!75!black] (0,0) -- (0,0.55) node[above left] {$\sigma_2\mathbf{v}_2$};
  \end{scope}
  \node at (8.4,-1.75) {image ellipse};

  \draw[map] (1.75,0.45) -- (2.8,0.45) node[midway, above] {$W^*$};
  \draw[map] (5.85,0.45) -- (6.95,0.45) node[midway, above] {$V$};
  \node[blue!70!black] at (4.1,1.55) {$\Sigma$};
\end{tikzpicture}

  \caption{The factors $W^*$, $\Sigma$, and $V$ rotate, stretch, and reorient the unit ball.}
  \label{fig:matrix-svd-unit-ball-to-ellipse}
\end{illustrationfigure}

\begin{perspective} \label{per:operator-to-matrix-svd-bridge} \hlt{Operator-to-Matrix SVD Bridge}\\
In the bases from Theorem \ref{thm:rectangular-operator-singular-vector-decomposition}, the matrix of $T$ is rectangular diagonal with positive entries $s_1,\ldots,s_r$ and zeros elsewhere.\\
Changing to standard orthonormal bases produces the matrix SVD in Theorem \ref{thm:matrix-singular-value-decomposition} without changing the singular data.
\end{perspective}

\begin{proposition} \label{prop:basic-singular-value-identities} \hlt{Basic Singular Value Identities}\\
Let $A\in\operatorname{Mat}(m,n,\mathbb{F})$ have singular values $\sigma_1\geq\cdots\geq\sigma_q\geq0$, where $q\defeq\min(m,n)$ and $r\defeq\operatorname{rank}(A)$. Let $c\in\mathbb{F}$.
\begin{enumerate}[label=(\roman*)]
\setlength{\itemsep}{0pt}
\item The positive eigenvalues of $A^*A$ and $AA^*$ are $\sigma_1^2,\ldots,\sigma_r^2$.
\item $\sum_{j=1}^{q}\sigma_j^2=\operatorname{tr}(A^*A)=\operatorname{tr}(AA^*)=\norm{A}_F^2$, where $\norm{A}_F^2\defeq\operatorname{tr}(A^*A)$.
\item $A$ and $A^*$ have the same singular values. Over $\mathbb{C}$, so do $A^t$ and $\overline{A}$.
\item The singular values of $cA$ are $\abs{c}\sigma_1,\ldots,\abs{c}\sigma_q$.
\end{enumerate}
\begin{proof}
Fix an SVD $A=V\Sigma W^*$ by Theorem \ref{thm:matrix-singular-value-decomposition}. Then $A^*A=W(\Sigma^*\Sigma)W^*$ and $AA^*=V(\Sigma\Sigma^*)V^*$.
\begin{enumerate}[label=(\roman*)]
\setlength{\itemsep}{0pt}
\item The positive eigenvalues of both Gram matrices are $\sigma_1^2,\ldots,\sigma_r^2$ by the two unitary similarities and Theorem \ref{thm:similarity-preserves-characteristic-polynomial}(i)--(ii).
\item $\operatorname{tr}(A^*A)=\operatorname{tr}(AA^*)=\sum\limits_{j=1}^{q}\sigma_j^2=\norm{A}_F^2$ by Theorem \ref{thm:similarity-preserves-trace-and-determinant}(i).
\item Claim for $A^*$ holds by (i). Over $\mathbb{C}$, $(A^t)^*A^t=(AA^*)^t$, $(\overline{A})^*\overline{A}=(A^*A)^t$, with unchanged Gram spectra by Proposition \ref{prop:transpose-conjugate-adjoint-characteristic-polynomial}.
\item The singular values of $cA$ are $\abs{c}\sigma_1,\ldots,\abs{c}\sigma_q$ by $(cA)^*(cA)=\abs{c}^2A^*A$.
\end{enumerate}
\hlt{Structural Commentary}: Unitary Gram reductions transfer spectra, traces, adjoints, and scalar scaling.\\
\hlt{Key Insight}: Both Gram spectra contain the squared entries of $\Sigma$.
\end{proof}
\end{proposition}

\begin{corollary} \label{cor:square-invertibility-singular-value-dictionary} \hlt{Square Invertibility and Singular Values}\\
Let $A\in\operatorname{Mat}(n,n,\mathbb{F})$, $\mathbb{F}\in\{\mathbb{R},\mathbb{C}\}$, have singular values $\sigma_1\geq\cdots\geq\sigma_n\geq0$. The following are equivalent:
\begin{enumerate}[label=(\roman*)]
\setlength{\itemsep}{0pt}
\item $A$ is invertible;
\item $\operatorname{rank}(A)=n$;
\item $\sigma_n>0$;
\item $A^*A\succ0$.
\end{enumerate}
\begin{proof}
\begin{enumerate}[label=(\roman*)]
\setlength{\itemsep}{0pt}
\item[$(i)$] $\Rightarrow (ii)$ $\operatorname{rank}(A)=n$ by Corollary \ref{cor:square-matrix-invertible-iff-full-rank}.
\item[$(ii)$] $\Rightarrow (iii)$ $\sigma_n>0$ by the definition of the full singular-value list.
\item[$(iii)$] $\Rightarrow (iv)$ Since all $n$ singular values are positive, $\operatorname{rank}(A)=n$ and $A^*A\succ0$ by Proposition \ref{prop:rectangular-gram-products-are-psd}(v).
\item[$(iv)$] $\Rightarrow (i)$ $\ker(A)=\{\mathbf{0}\}$ by $(\forall\mathbf{x}\in\mathbb{F}^n)\,[A\mathbf{x}=\mathbf{0}\Rightarrow\mathbf{x}^*A^*A\mathbf{x}=0\Rightarrow\mathbf{x}=\mathbf{0}]$.\\
Thus $\operatorname{rank}(A)=n$ by Theorem \ref{thm:rank-nullity}, and $A$ is invertible by Corollary \ref{cor:square-matrix-invertible-iff-full-rank}.
\end{enumerate}
\hlt{Structural Commentary}: Rank links invertibility and $\sigma_n$; the Gram form detects injectivity.\\
\hlt{Key Insight}: $\sigma_n>0$ is exactly full rank for a square matrix.
\end{proof}
\end{corollary}

\begin{definition} \label{def:compact-singular-value-decomposition} \hlt{Compact Singular Value Decomposition}\\
Let $A\in\operatorname{Mat}(m,n,\mathbb{F})$ have rank $r\geq1$. Factorisation $A=X\Sigma_rY^*$ is a \hlt{compact singular value decomposition} if $X\in\operatorname{Mat}(m,r,\mathbb{F})$, $Y\in\operatorname{Mat}(n,r,\mathbb{F})$, $X^*X=Y^*Y=I_r$, and $\Sigma_r=\operatorname{diag}(\sigma_1,\ldots,\sigma_r)$ with $\sigma_1\geq\cdots\geq\sigma_r>0$.
\end{definition}

\begin{theorem} \label{thm:compact-singular-value-decomposition} \hlt{Compact Singular Value Decomposition}\\
Let $A\in\operatorname{Mat}(m,n,\mathbb{F})$ have rank $r\geq1$.
\begin{enumerate}[label=(\roman*)]
\setlength{\itemsep}{0pt}
\item If $A=V\Sigma W^*$ is an SVD and $X,Y$ are the first $r$ columns of $V,W$, then $A=X\Sigma_rY^*$.
\item Conversely, if $X\in\operatorname{Mat}(m,r,\mathbb{F})$ and $Y\in\operatorname{Mat}(n,r,\mathbb{F})$ have orthonormal columns and $\Sigma_r=\operatorname{diag}(\sigma_1,\ldots,\sigma_r)$,\\
with $\sigma_1\geq\cdots\geq\sigma_r>0$, then $X\Sigma_rY^*$ has positive singular values $\sigma_1,\ldots,\sigma_r$.
\end{enumerate}
\begin{proof}
\begin{enumerate}[label=(\roman*)]
\setlength{\itemsep}{0pt}
\item For $V=[X\ X_0]$ and $W=[Y\ Y_0]$, $A=V\Sigma W^*=X\Sigma_rY^*$ by block multiplication.
\item Set $B\defeq X\Sigma_rY^*$. Then $B^*B=Y\Sigma_r^2Y^*$ by $X^*X=I_r$. Complete $Y$ to a unitary $\widetilde Y=[Y\ Y_0]$ by Theorem \ref{thm:extend-isometric-columns-to-unitary-matrix}(ii).\\
Then $\widetilde Y^*B^*B\widetilde Y=\operatorname{diag}(\Sigma_r^2,0)$, so $B$ has positive singular values $\sigma_1,\ldots,\sigma_r$ by Corollary \ref{cor:positive-semidefinite-square-root-formula}(ii).
\end{enumerate}
\hlt{Structural Commentary}: Block truncation constructs the compact form; unitary completion recovers its Gram spectrum.\\
\hlt{Key Insight}: $B^*B=Y\Sigma_r^2Y^*$ retains exactly the positive singular data.
\end{proof}
\end{theorem}

\begin{theorem} \label{thm:compact-svd-subspaces} \hlt{Column and Kernel Spaces from Compact SVD}\\
Let $A=X\Sigma_rY^*$ be a compact SVD. Then:
\begin{enumerate}[label=(\roman*)]
\setlength{\itemsep}{0pt}
\item $\operatorname{col}(A)=\operatorname{col}(X)=(\ker A^*)^{\perp}$;
\item $\ker(A)=\operatorname{col}(Y)^{\perp}=(\operatorname{col}(A^*))^{\perp}$.
\end{enumerate}
\begin{proof}
\begin{enumerate}[label=(\roman*)]
\setlength{\itemsep}{0pt}
\item $\operatorname{col}(A)=\operatorname{col}(X)$ by $A=X\Sigma_rY^*$ and $AY\Sigma_r^{-1}=X$, while $\operatorname{col}(A)=(\ker A^*)^\perp$ by Proposition \ref{prop:adjoint-kernel-range-orthogonality}(iv).
\item $(\forall\mathbf{z}\in\mathbb{F}^n)\,A\mathbf{z}=\mathbf{0}\Longleftrightarrow Y^*\mathbf{z}=\mathbf{0}\Longleftrightarrow\mathbf{z}\perp\operatorname{col}(Y)$ by $X^*X=I_r$ and invertibility of $\Sigma_r$. $\ker(A)=(\operatorname{col}(A^*))^\perp$ by Proposition \ref{prop:adjoint-kernel-range-orthogonality}(ii).
\end{enumerate}
\hlt{Structural Commentary}: Invertible middle scaling transfers the image and isolates the right orthogonal kernel.\\
\hlt{Key Insight}: $X\Sigma_r$ is injective on the right singular coordinates.
\end{proof}
\end{theorem}

Figure \ref{fig:compact-svd-subspaces} summarizes the active right singular subspace, its image, and the kernel annihilated by $A$.

\begin{illustrationfigure}[H]
  \centering
  \begin{tikzpicture}[
  >=Stealth,
  x=0.86cm,
  y=0.88cm,
  every node/.style={font=\small},
  box/.style={draw=gray!65, rounded corners=3pt, minimum width=2.05cm, minimum height=0.58cm, align=center},
  map/.style={->, line width=0.9pt, blue!70!black},
  kill/.style={->, line width=0.9pt, gray!70}
]
  \node[box] (active) at (0,0.7) {$\operatorname{col}(Y)$};
  \node[box] (kernel) at (0,-0.7) {$\ker(A)$};
  \node[box] (image) at (4.2,0.7) {$\operatorname{col}(A)=\operatorname{col}(X)$};
  \node[box] (zero) at (4.2,-0.7) {$\{\mathbf{0}\}$};

  \draw[map] (active) -- (image)
    node[midway, above] {$A$};
  \draw[kill] (kernel) -- (zero)
    node[midway, below] {$A$};
\end{tikzpicture}

  \caption{Compact SVD: $A$ maps $\operatorname{col}(Y)$ onto $\operatorname{col}(X)$ and kills $\ker(A)=\operatorname{col}(Y)^\perp$.}
  \label{fig:compact-svd-subspaces}
\end{illustrationfigure}

\begin{theorem} \label{thm:compact-svd-unitary-ambiguity} \hlt{Uniqueness of Compact SVD up to Singular Subspaces}\\
Let $A=X_1\Sigma_rY_1^*=X_2\Sigma_rY_2^*$ be compact SVDs with the same $\Sigma_r$.\\
Then $\exists$ unitary matrix $U\in\operatorname{Mat}(r,r,\mathbb{F})$, block diagonal on equal-value blocks of $\Sigma_r$, such that $X_1=X_2U$, $Y_1=Y_2U$, $U\Sigma_r=\Sigma_rU$.
\begin{proof}
\begin{enumerate}[label=(\roman*)]
\setlength{\itemsep}{0pt}
\item \hlt{Coordinate Changes}: The equalities $\operatorname{col}(X_1)=\operatorname{col}(X_2)$ and $\operatorname{col}(Y_1)=\operatorname{col}(Y_2)$ hold by Theorem \ref{thm:compact-svd-subspaces}.\\
For $U\defeq X_2^*X_1$ and $Z\defeq Y_2^*Y_1$, $X_1=X_2U$, $Y_1=Y_2Z$, and $U,Z$ are unitary by the common column spaces and orthonormality.
\item \hlt{Gram Coupling}: $\Sigma_r=U\Sigma_rZ^*$ by substitution, so $U\Sigma_r=\Sigma_rZ$. $U\Sigma_r^2U^*=\Sigma_r^2$ by this equality and its adjoint.
\item \hlt{Block Ambiguity}: The matrix $U$ is block diagonal on the equal-value blocks of $\Sigma_r$ by $U\Sigma_r^2=\Sigma_r^2U$ and Lemma \ref{lem:block-diagonal-commutant-distinct-scalars}(ii).\\
$U\Sigma_r=\Sigma_rU$ by this block structure.
\item \hlt{Common Change}: $Z=U$ by $U\Sigma_r=\Sigma_rZ=\Sigma_rU$ and invertibility of $\Sigma_r$.
\end{enumerate}
\hlt{Structural Commentary}: Common subspaces create two unitary changes; Gram coupling identifies them blockwise.\\
\hlt{Key Insight}: Commutation with $\Sigma_r^2$ confines motion to equal-value blocks.
\end{proof}
\end{theorem}

\begin{proposition} \label{prop:matrix-modulus-from-svd} \hlt{Modulus from SVD}\\
Let $A=V\Sigma W^*$ be an SVD. Then:
\begin{enumerate}[label=(\roman*)]
\setlength{\itemsep}{0pt}
\item $\abs{A}=W(\Sigma^*\Sigma)^{1/2}W^*$;
\item $\abs{A^*}=V(\Sigma\Sigma^*)^{1/2}V^*$.
\end{enumerate}
\begin{proof}
Set $D_R\defeq(\Sigma^*\Sigma)^{1/2}$ and $D_L\defeq(\Sigma\Sigma^*)^{1/2}$.
\begin{enumerate}[label=(\roman*)]
\setlength{\itemsep}{0pt}
\item The matrix $WD_RW^*$ is positive semidefinite by Theorem \ref{thm:positive-semidefinite-stability-rules}(ii) and squares to $A^*A$. Hence $\abs{A}=WD_RW^*$ by Corollary \ref{cor:positive-semidefinite-square-root-formula}(i).
\item The matrix $VD_LV^*$ is positive semidefinite by Theorem \ref{thm:positive-semidefinite-stability-rules}(ii) and squares to $AA^*$. Hence $\abs{A^*}=VD_LV^*$ by Corollary \ref{cor:positive-semidefinite-square-root-formula}(i).
\end{enumerate}
\hlt{Structural Commentary}: Unitary conjugation transports each diagonal root to its Gram matrix.\\
\hlt{Key Insight}: The two moduli express $\Sigma$ in opposite singular bases.
\end{proof}
\end{proposition}

\begin{theorem} \label{thm:rectangular-polar-decomposition} \hlt{Rectangular Polar Decomposition}\\
Let $A=V\Sigma W^*\in\operatorname{Mat}(m,n,\mathbb{F})$ be an SVD.
\begin{enumerate}[label=(\roman*)]
\setlength{\itemsep}{0pt}
\item If $m\geq n$, let $V_n$ be the first $n$ columns of $V$ and set $U=V_nW^*$. Then $U^*U=I_n$ and $A=U\abs{A}$.
\item If $m\leq n$, let $W_m$ be the first $m$ columns of $W$ and set $U=VW_m^*$. Then $UU^*=I_m$ and $A=\abs{A^*}U$.
\item If $m=n$, then $U=VW^*$ is unitary and $A=U\abs{A}=\abs{A^*}U$.
\end{enumerate}
\begin{proof}
\begin{enumerate}[label=(\roman*)]
\setlength{\itemsep}{0pt}
\item For $D_R\defeq(\Sigma^*\Sigma)^{1/2}$, $V_nD_R=V\Sigma$; hence $U\abs{A}=A$ by Proposition \ref{prop:matrix-modulus-from-svd}(i), while $U^*U=WV_n^*V_nW^*=I_n$.
\item By (i) with $A^*=W\Sigma^*V^*$, $A^*=U_0\abs{A^*}$ for $U_0\defeq W_mV^*$ and $U_0^*U_0=I_m$.\\
For $U=U_0^*=VW_m^*$, $A=\abs{A^*}U$ and $UU^*=I_m$ by taking adjoints.
\item If $m=n$, then $V_n=V$, $W_m=W$, and $U=VW^*$ is unitary; both identities hold by (i)--(ii).
\end{enumerate}
\hlt{Structural Commentary}: Each dimensional regime absorbs the compatible singular bases into a one-sided isometry.\\
\hlt{Key Insight}: $VW^*$ preserves the diagonal stretch encoded by the modulus.
\end{proof}
\end{theorem}

\begin{computationmethod} \label{meth:svd-polar-computation} \hlt{SVD and Polar Factor Computation}\\
Given $A\in\operatorname{Mat}(m,n,\mathbb{F})$, method constructs full SVD, its compact form when $A\neq0$, and the dimension-compatible polar factorisation:
\begin{enumerate}[label=(\roman*)]
\setlength{\itemsep}{0pt}
\item Set $r\defeq\operatorname{rank}(A)$ and diagonalise $A^*A$ with a unitary eigenvector matrix $W=[\mathbf{w}_1\ \cdots\ \mathbf{w}_n]$.\\
For $r\geq1$, order its positive eigenvalues as $\sigma_1^2\geq\cdots\geq\sigma_r^2>0$; for $r=0$, the positive list is empty.
\item Set $(\forall j\in J_r)\,\mathbf{v}_j\defeq\sigma_j^{-1}A\mathbf{w}_j$, and complete $(\mathbf{v}_1,\ldots,\mathbf{v}_r)$ to an orthonormal basis of $\mathbb{F}^m$. Let $V$ contain this basis.
\item Form the rectangular diagonal $\Sigma$ with diagonal $\sigma_1,\ldots,\sigma_r,0,\ldots,0$. Verify $AW=V\Sigma$, and hence $A=V\Sigma W^*$.
\item If $r\geq1$, set $X\defeq[\mathbf{v}_1\ \cdots\ \mathbf{v}_r]$, $Y\defeq[\mathbf{w}_1\ \cdots\ \mathbf{w}_r]$, and $\Sigma_r\defeq\operatorname{diag}(\sigma_1,\ldots,\sigma_r)$. Then $A=X\Sigma_rY^*$ by Theorem \ref{thm:compact-singular-value-decomposition}(i).
\item Compute the moduli by Proposition \ref{prop:matrix-modulus-from-svd}. For $m\geq n$, set $U\defeq V_nW^*$ and obtain $A=U\abs{A}$.\\
For $m\leq n$, set $U\defeq VW_m^*$ and obtain $A=\abs{A^*}U$ by Theorem \ref{thm:rectangular-polar-decomposition}.
\end{enumerate}
\end{computationmethod}

\begin{theorem} \label{thm:polar-factor-uniqueness} \hlt{Uniqueness of the Positive Polar Factor}\\
Let $A\in\operatorname{Mat}(m,n,\mathbb{F})$.
\begin{enumerate}[label=(\roman*)]
\setlength{\itemsep}{0pt}
\item If $m\geq n$ and $A=UB$ with $U^*U=I_n$ and $B\succeq0$, then $B=\abs{A}$. If $\operatorname{rank}(A)=n$, then $U=A\abs{A}^{-1}$ is unique.
\item If $m\leq n$ and $A=CU$ with $UU^*=I_m$ and $C\succeq0$, then $C=\abs{A^*}$. If $\operatorname{rank}(A)=m$, then $U=\abs{A^*}^{-1}A$ is unique.
\end{enumerate}
\begin{proof}
\begin{enumerate}[label=(\roman*)]
\setlength{\itemsep}{0pt}
\item From $A^*A=BU^*UB=B^2$, $B=\abs{A}$ by Corollary \ref{cor:positive-semidefinite-square-root-formula}(i). If $\operatorname{rank}(A)=n$, then $A^*A\succ0$ by Proposition \ref{prop:rectangular-gram-products-are-psd}(v).\\
Thus $\ker\abs{A}=\{\mathbf{0}\}$ and $U=A\abs{A}^{-1}$ by Corollary \ref{cor:square-matrix-invertible-iff-full-rank}.
\item From $AA^*=CUU^*C=C^2$, $C=\abs{A^*}$ by Corollary \ref{cor:positive-semidefinite-square-root-formula}(i). If $\operatorname{rank}(A)=m$, then $AA^*\succ0$ by Proposition \ref{prop:rectangular-gram-products-are-psd}(vi).\\
Thus $\ker\abs{A^*}=\{\mathbf{0}\}$ and $U=\abs{A^*}^{-1}A$ by Corollary \ref{cor:square-matrix-invertible-iff-full-rank}.
\end{enumerate}
\hlt{Structural Commentary}: Gram squaring fixes each positive factor; full rank makes it invertible.\\
\hlt{Key Insight}: An invertible modulus uniquely isolates the one-sided isometry.
\end{proof}
\end{theorem}

\begin{theorem} \label{thm:normal-iff-polar-factors-commute} \hlt{Normality and Polar Commutation}\\
Let $A\in\operatorname{Mat}(n,n,\mathbb{F})$ and suppose $A=U\abs{A}$ with $U$ unitary. Then $A$ is normal if and only if $U\abs{A}=\abs{A}U$.
\begin{proof}
\begin{enumerate}[label=\roman*.]
\setlength{\itemsep}{0pt}
\item[$(\Rightarrow)$] $\abs{A}^2=A^*A=AA^*=U\abs{A}^2U^*$ by normality. $U\abs{A}U^*$ is positive semidefinite by Theorem \ref{thm:positive-semidefinite-stability-rules}(ii), and $(U\abs{A}U^*)^2=\abs{A}^2$.\\
Thus $U\abs{A}U^*=\abs{A}$ by Corollary \ref{cor:positive-semidefinite-square-root-formula}(i), and $U\abs{A}=\abs{A}U$.
\item[$(\Leftarrow)$] $AA^*=U\abs{A}^2U^*=\abs{A}^2=A^*A$ by $U\abs{A}=\abs{A}U$.
\end{enumerate}
\hlt{Structural Commentary}: Normality identifies positive roots; commutation reverses the Gram identity.\\
\hlt{Key Insight}: Root uniqueness converts Gram equality into polar commutation.
\end{proof}
\end{theorem}

\begin{perspective} \label{per:polar-decomposition-spectral-analogy} \hlt{Polar Decomposition Separates Stretch and Direction}\\
In square factorisation $A=U\abs{A}$, positive factor records metric stretch, unitary factor records direction, paralleling scalar modulus and phase. For normal $A$, these factors commute by Theorem \ref{thm:normal-iff-polar-factors-commute}.\\
In rectangular case, Theorem \ref{thm:rectangular-polar-decomposition} supplies only dimension-compatible one-sided isometry.
\end{perspective}

\begin{proposition} \label{prop:absolute-determinant-positive-polar-factor} \hlt{Absolute Determinant Comes from the Positive Polar Factor}\\
Let $\mathcal{V}\neq\{\mathbf{0}\}$ and $T\in\LL(\mathcal{V})$. Then $\abs{\det T}=\det\sqrt{T^*T}$.
\begin{proof}
Set $P\defeq\sqrt{T^*T}$ and fix an isometry $S$ satisfying $T=SP$ by Theorem \ref{thm:polar-decomposition}. Then $\abs{\det S}=1$ by Proposition \ref{prop:isometry-determinant-absolute-one}(ii).\\
Fix an orthonormal eigenbasis of $P$ by Theorem \ref{thm:real-spectral-theorem} when $\mathbb{F}=\mathbb{R}$ and Theorem \ref{thm:complex-spectral-theorem} when $\mathbb{F}=\mathbb{C}$.\\
Its eigenvalues are nonnegative by Theorem \ref{thm:positive-operator-characterization}(i)$\Rightarrow$(ii), so $\det P\geq0$ by Theorem \ref{thm:determinant-basic-laws}(v).\\
Thus $\abs{\det T}=\abs{\det S}\abs{\det P}=\det P$ by Theorem \ref{thm:determinant-basic-laws}(vi).\\
\hlt{Structural Commentary}: Polar multiplicativity separates unit-modulus motion from the nonnegative spectral product.\\
\hlt{Key Insight}: The positive factor alone carries the absolute determinant.
\end{proof}
\end{proposition}
